% =====================================================================
% 4.1 Anyagok es modszerek -- In vivo kiserletek (Kereg--GPe anatomia)
% Forras: KOKI laborjegyzokonyv 2026 nyar (GPE16-GPE24).
% FIGYELEM: a [?] jelolesek a jegyzokonyvben hianyzo/ellenorizendo adatok.
% =====================================================================

\subsection{In vivo kísérletek (Kéreg--GPe anatómia)}

\subsubsection{Kísérleti állatok és sebészeti beavatkozás}

A kéreg--GPe kapcsolat kísérletes feltérképezéséhez eddig összesen hét
egérben (GPE16--GPE24) végeztünk retrográd pályajelölést; a beavatkozásokat
a kutatócsoport belső laborjegyzőkönyvében dokumentáljuk. Az állatokat
izoflurán altatásban, sztereotaxikus készülékben rögzítettük, majd a GPe
eltérő antero-posterior, medio-laterális és dorzo-ventrális koordinátáira
kis térfogatú Fluorogold retrográd jelölőanyagot juttattunk be
(\ref{tab:mutetek}. táblázat). A jelölőanyagot 12--19~\textmu m
hegyátmérőjű üvegpipettán keresztül, iontoforetikusan adtuk be
(0,5--1~\textmu A, 2~s be / 2~s ki ciklusokban, 10 percen át; a GPE17
állat esetében 2 percig). A retrográd jelölőanyagot a GPe-be vetítő sejtek
axonterminálisai veszik fel, és a sejttestükbe transzportálják, így az
injekció helyétől függően azonosítható, mely kérgi területek küldenek oda
bemenetet.

\begin{table}[ht]
\centering
\caption{Az eddig elvégzett Fluorogold-injekciók sztereotaxikus koordinátái
(bregmához viszonyítva, mm) és kimenetele.}
\label{tab:mutetek}
\begin{tabular}{lrrrll}
\hline
Állat & AP & ML & DV & Iontoforézis & Kimenetel \\
\hline
GPE16 & $-0{,}22$ & $1{,}9$  & $-3{,}4$ & 0,5~\textmu A, 10~p & kiértékelve \\
GPE17 & $-1{,}0$  & $2{,}3$  & $-3{,}2$ & 0,5~\textmu A, 2~p  & injekció rossz helyre került \\
GPE19 & [?]       & [?]      & [?]      & [?]                 & metszési hiba \\
GPE20 & $-1{,}0$  & $2{,}3$  & $-3{,}2$ & 0,5~\textmu A, 10~p & kis injekciós hely \\
GPE22 & $-0{,}2$  & $1{,}8$  & $-3{,}4$ & 1~\textmu A, 10~p   & részleges kéreg \\
GPE23 & $-1{,}0$  & $-2{,}4$ & $-3{,}4$ & 1~\textmu A, 10~p   & részleges kéreg \\
GPE24 & $-0{,}5$  & $-2{,}4$ & $-3{,}4$ & 1~\textmu A, 10~p   & kiértékelve \\
\hline
\end{tabular}
\end{table}

A hét beavatkozásból kettő (GPE16, GPE24) jutott el a teljes kiértékelésig.
A többi esetben technikai okok akadályozták a feldolgozást: a GPE17 állatnál
a bregma téves azonosítása miatt az injekció a céltérfogaton kívülre került;
a GPE19 esetében a nem kellően stabil agy miatt a metszés éppen a releváns
területen sérült; a GPE20 állatnál a pipetta eltömődése miatt háromszor
kellett újrapozicionálni, és az így létrejött injekciós hely túl kicsi lett;
a GPE22 és GPE23 állatoknál a perfúzió tökéletlensége miatt (a szív hibás
felnyitása, utólagos posztfixálás) a kéreg nem maradt teljes egészében
feldolgozható állapotban. Ezek a technikai kudarcok a módszer tanulási
görbéjének részét képezik: a legutolsó, GPE24 jelű beavatkozásnál mind az
iontoforézis árama, mind a perfúzió és a metszés hibátlanul zajlott, ami azt
jelzi, hogy a protokoll mostanra stabilizálódott.

\subsubsection{Perfúzió és szövettani feldolgozás}

Az injekciót követő 6--12. napon \emph{[ellenőrizendő: a jegyzőkönyvben öt műtét- és öt perfúziódátum szerepel hét állatra]} az állatokat mély altatásban,
transzkardiálisan perfundáltuk 4\%-os paraformaldehid-oldattal
(30--60 perc, kb.\ 200~ml). A GPE20--GPE24 állatok esetében a fixálást
megelőzően glutáraldehid-tartalmú oldatot is alkalmaztunk. Az agyakból
50~\textmu m vastag metszeteket készítettünk, amelyeket két párhuzamos
sorozatba (5 kérgi + 5 talamikus üveg) gyűjtöttünk.

A Fluorogold immunhisztokémiai kimutatását két napos protokollal végeztük.
Első nap: endogén peroxidáz-blokk (30-szorosan hígított hidrogén-peroxid),
majd blokkolás 0,2\% Triton X-100 és 1\% normál szamárszérum (NDS)
elegyében -- a GPE20 állat esetében NDS helyett BSA-val --, ezt követően
nyúlban termelt poliklonális anti-Fluorogold primer antitest
(Chemicon, 1:30\,000) inkubálása. Második nap: biotinilált anti-nyúl
szekunder antitest (1:500), avidin--biotin komplex (ABC), majd
nikkel-intenzifikált diaminobenzidin (DAB-Ni) kromogén reakció
(8~ml PB + 0,5~ml DAB + 0,0033~g ammónium-klorid + 400~\textmu l
nikkel-ammónium-szulfát; hívás 5~ml desztillált vízben 5~\textmu l 30\%-os
H\textsubscript{2}O\textsubscript{2}-vel). A metszeteket
krómzselatinos tárgylemezre húztuk fel és DPX-szel fedtük.
A GPE22--GPE24 állatok második metszetsorozatát fluoreszcens kimutatásra
tartottuk fenn: ezeket foszfátpufferből húztuk fel és Vectashield
fedőanyaggal fedtük.

\subsubsection{Képalkotás és kiértékelés}

A DAB-Ni-vel festett metszeteket slide scannerrel áttekintő felbontásban
digitalizáltuk, a részletes vizsgálathoz pedig konfokális mikroszkópiával
készítettünk nagyfelbontású felvételeket. A képeken azonosítjuk a GPe-be
vetítő, retrográd módon megjelölt kérgi sejttesteket, és rögzítjük az
injekció pontos helyét és kiterjedését.

\emph{[Ide kerül még: a sejtszámlálás és a térképre illesztés (regisztráció
egy referencia-atlaszhoz) módszertana és statisztikai kiértékelése, valamint
az injekciós koordináták összevetése az Abecassis és mtsai.\ (2020)
1.~táblázatában közölt GPe-koordinátákkal mint irodalmi referenciaponttal.]}

\subsubsection{Tervezett anterográd validáció}

Kiegészítő, a projekt második felében tervezett módszerként Rbp4-Cre
transzgén egerek agykérgébe SynaptoTag vírust injektálunk, amely szelektíven
jelöli meg az 5. rétegi (L5) sejtek szinaptikus terminálisait; a hasonló,
kettős fluoreszcens SynaptoTag-konstrukciókat korábban sikeresen alkalmazták
projekció-specifikus terminálisok anatómiai azonosítására \cite{xia2020}.
A virális genom kizárólag L5 kérgi sejtekben aktiválódik, a virális fehérje
pedig a sejtek idegvégződéseibe transzportálódik. Ezzel a megközelítéssel
anterográd irányból, közvetlenül ellenőrizhető és megerősíthető a retrográd
(Fluorogold-alapú) kísérletekből származó topográfiai térkép: a GPe-ből
retrográd módon megjelölődött kérgi régiók esetén az anterográd
SynaptoTag-jelölésnek is meg kell jelennie a GPe-ben, míg a retrográd jelet
nem mutató kérgi régiók esetén ennek hiányát várjuk.
